\chapter{Quantum error correction}
\label{ch:qec}

The previous chapter treated qubits and gates as perfect. They are not.
Part~II showed that a real qubit relaxes on a timescale $T_1$
(\cref{stmt:t1}), dephases on $T_\phi$ (\cref{stmt:tphi}), and that its
evolution is the noisy, irreversible dynamics of a Lindblad master
equation (\cref{stmt:lindblad}) rather than a clean unitary; Part~IV
catalogued the hardware channels that cause it. A computation of any
useful length applies thousands to billions of gates, and at realistic
per-gate error rates the logical state is scrambled long before the
answer appears. Classical computers face the same threat and defeat it
with redundancy --- store each bit in triplicate and take a majority
vote. This chapter asks whether the same can be done for quantum
information, and the answer is the most surprising constructive result
in the subject: yes, despite three obstacles that look fatal --- the
no-cloning theorem (\cref{thm:no-cloning}) forbids copying, quantum
errors form a continuum rather than a discrete set, and any measurement
that checks for an error threatens to collapse the very superposition we
are protecting. We will see all three dissolved at once. The vehicle is
the stabiliser formalism of \cref{stmt:gottesman-knill} --- the same
bookkeeping that made Clifford circuits classically simulable, here
turned into a shield --- developed from the three-qubit repetition code
up through the Shor and Steane codes, the surface code that leads
present hardware, and finally the bosonic codes that protect a logical
qubit inside a single superconducting cavity, closing the loop back to
Part~IV.

% --------------------------------------------------------------
\section{Why error correction is hard --- and possible}
\label{sec:why-qec}

\begin{phenomenon}
Take one qubit in the state $\alpha\ket{0}+\beta\ket{1}$ and leave it in
a real device. Energy relaxation drives it towards $\ket{0}$; random
frequency fluctuations scramble the relative phase of $\alpha$ and
$\beta$. The classical fix --- copy the bit three times and majority
vote --- fails on contact with quantum mechanics. We cannot copy
$\alpha,\beta$ (no-cloning). The errors are not discrete flips but a
continuum: an over-rotation by a tiny angle, a partial phase slip, a
little entanglement with a stray mode. And to find out whether an error
occurred we would ordinarily measure the qubit, which collapses the
superposition we are trying to save. Every ingredient of the classical
recipe is forbidden. Remarkably, a different recipe works.
\end{phenomenon}

The escape has three matching moves. We encode without copying, by
spreading one logical qubit across several physical qubits in an
\emph{entangled} state whose amplitudes are still $\alpha,\beta$; we
measure without collapsing, by reading \emph{parity-check} operators
that reveal whether an error happened but not what the logical state is;
and we tame the continuum by the following theorem, which says that
correcting a handful of discrete errors automatically corrects them all.

\begin{statement}[\textbf{Error discretisation}]
\label{stmt:digitisation}
Model noise on a qubit as a quantum channel with Kraus operators
$\{\op A_k\}$ (\cref{stmt:kraus}). Each $\op A_k$ expands in the Pauli
basis $\op A_k=\sum_{P\in\{\id,\hat\sigma_x,\hat\sigma_y,\hat\sigma_z\}}
c_{kP}\,P$. If a code can detect and correct each of the discrete errors
$\hat\sigma_x,\hat\sigma_y,\hat\sigma_z$ on a given qubit --- and
identify which occurred --- then measuring the error syndrome projects
any continuous or coherent error on that qubit onto one definite Pauli
branch, which the same recovery then undoes. Correcting the discrete set
$\{\hat\sigma_x,\hat\sigma_y,\hat\sigma_z\}$ therefore corrects an
arbitrary single-qubit error, including entanglement with the
environment.
\end{statement}

\begin{proof}[Derivation]
Consider a coherent over-rotation of qubit $j$ by a small angle $\theta$,
\begin{equation}
\op R=\exp\!\bigl(-i\theta\,\hat\sigma_x^{(j)}\bigr)
=\cos\theta\,\id-i\sin\theta\,\hat\sigma_x^{(j)},
\end{equation}
acting on an encoded state $\ket{\psi_L}$ that lives in the $+1$
eigenspace of the code's parity checks. The post-error state
$\op R\ket{\psi_L}=\cos\theta\,\ket{\psi_L}-i\sin\theta\,
\hat\sigma_x^{(j)}\ket{\psi_L}$
is a superposition of the clean state and the state with a single
$\hat\sigma_x^{(j)}$ error. The two branches carry \emph{different}
syndromes: $\ket{\psi_L}$ gives the trivial syndrome,
$\hat\sigma_x^{(j)}\ket{\psi_L}$ a nontrivial one. Measuring the
syndrome therefore collapses the superposition --- with probability
$\cos^2\theta$ to the clean branch (syndrome trivial, nothing to do) and
with probability $\sin^2\theta$ to the error branch (syndrome flags
qubit $j$, apply $\hat\sigma_x^{(j)}$ to recover). The continuous
parameter $\theta$ has been digitised into a yes/no event. The same
argument applied to a general $\op A_k=\sum_P c_{kP}P$ shows the
syndrome measurement projects onto a definite Pauli error, so a code
correcting $\{\hat\sigma_x,\hat\sigma_y,\hat\sigma_z\}$ on each qubit
needs nothing more. Crucially the syndrome reveals only \emph{which
Pauli} occurred, never the values of $\alpha,\beta$: the parity checks
are built to commute with the logical operators, so the logical
information survives the measurement untouched.
\end{proof}

\begin{example}[A parity check that learns nothing about the logical bit]
\label{ex:parity-check}
Take two physical qubits and the operator
$\hat\sigma_z^{(1)}\hat\sigma_z^{(2)}$. Its eigenvalue is $+1$ on
$\ket{00}$ and $\ket{11}$, and $-1$ on $\ket{01}$ and $\ket{10}$: it
measures the \emph{parity} of the two bits, not their values. Applied to
a logical state $\alpha\ket{00}+\beta\ket{11}$, it returns $+1$ with
certainty and leaves the state unchanged --- it has confirmed ``no
disagreement'' while learning nothing about $\alpha$ or $\beta$. If
instead a bit flip has produced $\alpha\ket{10}+\beta\ket{01}$, the same
measurement returns $-1$, flagging the error, again without exposing the
logical amplitudes. This single observation is the seed of every code in
the chapter.
\end{example}

\paragraph{Transition.}
The two-qubit parity check already does something no classical bit
inspection can. Wiring three of them together gives the smallest genuine
error-correcting code, where the mechanism --- encode, extract syndrome,
recover --- can be seen in full.

% --------------------------------------------------------------
\section{Repetition codes}
\label{sec:repetition}

\begin{phenomenon}
Suppose for a moment that the only error is the bit flip
$\hat\sigma_x$. Classically we would store $0\to000$, $1\to111$ and
majority vote. The quantum version keeps the spirit --- threefold
redundancy --- but replaces both the copying and the voting with
entanglement and parity checks, so that no measurement ever reads the
protected amplitudes.
\end{phenomenon}

\begin{statement}[\textbf{The three-qubit repetition code}]
\label{stmt:repetition}
Encode $\alpha\ket{0}+\beta\ket{1}\;\mapsto\;
\alpha\ket{000}+\beta\ket{111}$, so that
$\ket{0_L}=\ket{000}$ and $\ket{1_L}=\ket{111}$. The code space is the
$+1$ eigenspace of the two stabiliser generators
\begin{equation}
g_1=\hat\sigma_z^{(1)}\hat\sigma_z^{(2)},
\qquad
g_2=\hat\sigma_z^{(2)}\hat\sigma_z^{(3)}.
\end{equation}
Measuring $(g_1,g_2)$ returns a two-bit \emph{syndrome} that uniquely
identifies a single bit-flip error and the qubit it struck, without
disturbing $\alpha,\beta$; applying $\hat\sigma_x$ to that qubit
restores the state. The logical operators are
$\op{Z}_L=\hat\sigma_z^{(1)}$ and
$\op{X}_L=\hat\sigma_x^{(1)}\hat\sigma_x^{(2)}\hat\sigma_x^{(3)}$.
\end{statement}

\begin{proof}[Derivation]
The encoding is two CNOTs from the data qubit onto two blank qubits,
\begin{equation}
(\alpha\ket{0}+\beta\ket{1})\ket{0}\ket{0}
\;\xrightarrow{\;\mathrm{CNOT}_{1\to2}\,\mathrm{CNOT}_{1\to3}\;}\;
\alpha\ket{000}+\beta\ket{111},
\end{equation}
which entangles rather than clones (\cref{thm:no-cloning}): the
amplitudes $\alpha,\beta$ are never duplicated, only spread. Now suppose
$\hat\sigma_x^{(2)}$ strikes, giving $\alpha\ket{010}+\beta\ket{101}$.
The check $g_1=\hat\sigma_z^{(1)}\hat\sigma_z^{(2)}$ compares qubits $1$
and $2$, which now disagree, returning $-1$; $g_2$ compares qubits $2$
and $3$, also disagreeing, returning $-1$. The full syndrome table is
\begin{equation}
\begin{aligned}
\text{no error}&:(+1,+1), &\quad
\hat\sigma_x^{(1)}&:(-1,+1),\\
\hat\sigma_x^{(2)}&:(-1,-1), &\quad
\hat\sigma_x^{(3)}&:(+1,-1),
\end{aligned}
\end{equation}
so the four outcomes of $(g_1,g_2)$ point at the four possibilities ---
no error, or a flip on qubit $1$, $2$, or $3$ --- and recovery applies
$\hat\sigma_x$ to the indicated qubit. The checks never collapse the
logical state because each commutes with both logical operators:
$g_1$ shares qubit $1$ with $\op{Z}_L=\hat\sigma_z^{(1)}$ (they are both
$\hat\sigma_z$, so commute) and shares qubits $1,2$ with
$\op{X}_L=\hat\sigma_x^{(1)}\hat\sigma_x^{(2)}\hat\sigma_x^{(3)}$ (two
anticommuting pairs, an even number, so commute). The syndrome thus
lives in a separate degree of freedom from the encoded qubit, exactly as
in \cref{ex:parity-check}.
\end{proof}

\begin{statement}[\textbf{The phase-flip code by Hadamard duality}]
\label{stmt:phase-flip}
The bit-flip code is blind to $\hat\sigma_z$ errors: a phase flip on any
qubit commutes with $g_1,g_2$ and goes undetected. Conjugating the whole
construction by Hadamards on every qubit exchanges
$\hat\sigma_x\leftrightarrow\hat\sigma_z$
(\cref{eq:clifford-conjugation}) and produces the \emph{phase-flip}
code: $\ket{0_L}=\ket{+\!+\!+}$, $\ket{1_L}=\ket{-\!-\!-}$, with
stabilisers $\hat\sigma_x^{(1)}\hat\sigma_x^{(2)}$ and
$\hat\sigma_x^{(2)}\hat\sigma_x^{(3)}$, correcting one $\hat\sigma_z$
error.
\end{statement}

\begin{example}[Why one repetition code is not enough]
\label{ex:rep-insufficient}
Each repetition code corrects bit flips \emph{or} phase flips, never
both: the bit-flip code's distance against $\hat\sigma_z$ is $1$ (a
single $\hat\sigma_z^{(1)}$ is an undetected logical operation). Real
noise contains both, and a $\hat\sigma_y=i\hat\sigma_x\hat\sigma_z$ event
is simultaneously a bit and phase flip. The way forward is to nest a
phase-flip code outside a bit-flip code so that each kind of error is
caught by its own layer --- precisely Shor's nine-qubit construction
(\cref{sec:css}). Before building it, we abstract the pattern ---
``define the code by the operators that stabilise it'' --- into a
formalism that handles all codes at once.
\end{example}

\paragraph{Transition.}
Every feature of the repetition code --- a code space fixed by commuting
parity operators, a syndrome from their eigenvalues, logical operators
that commute with the checks, and a distance counting undetectable
errors --- generalises verbatim. That generalisation is the stabiliser
formalism.

% --------------------------------------------------------------
\section{The stabiliser formalism}
\label{sec:stabiliser}

\begin{phenomenon}
The repetition code was specified not by writing out its two code words
but, more economically, by the two operators that fix them. This is the
general lesson: a code space of dimension $2^k$ inside $2^n$ physical
qubits is pinned down by $n-k$ commuting Pauli operators, and everything
about error detection follows from how candidate errors commute or
anticommute with them. The same algebra that made Clifford circuits
simulable (\cref{stmt:gottesman-knill}) now classifies codes.
\end{phenomenon}

\begin{statement}[\textbf{Stabiliser codes}~\cite{gottesman1997}]
\label{stmt:stabiliser-code}
An $[[n,k,d]]$ stabiliser code encodes $k$ logical qubits in $n$
physical qubits. Its code space is the simultaneous $+1$ eigenspace of
an abelian subgroup $\mathcal S$ of the $n$-qubit Pauli group with
$-\id\notin\mathcal S$, generated by $n-k$ independent commuting Pauli
operators $g_1,\dots,g_{n-k}$ (the \emph{stabilisers}). A Pauli error
$E$ produces the syndrome $s_i\in\{0,1\}$ with $s_i=0$ if $[E,g_i]=0$
and $s_i=1$ if $\{E,g_i\}=0$. The \emph{logical operators} are the
Paulis that commute with all $g_i$ but lie outside $\mathcal S$ (the
normaliser quotient $N(\mathcal S)/\mathcal S$); the code \emph{distance}
$d$ is the minimum weight of such a logical operator. The code corrects
any error of weight up to $\lfloor(d-1)/2\rfloor$.
\end{statement}

\begin{proof}[Derivation]
Let $\ket{\psi}$ be in the code space, so $g_i\ket{\psi}=\ket{\psi}$ for
all $i$. If an error $E$ anticommutes with $g_i$, then
\begin{equation}
g_i\,E\ket{\psi}=-E\,g_i\ket{\psi}=-E\ket{\psi},
\end{equation}
so $E\ket{\psi}$ sits in the $-1$ eigenspace of $g_i$: measuring $g_i$
returns $-1$ and flags the error. The collection of these signs is the
syndrome, and because every $g_i$ commutes with every logical operator,
the measurement reveals the error without touching the encoded state.
Two errors $E,E'$ are confused only if they share a syndrome, i.e.\
$E'E^\dagger$ commutes with all $g_i$ and so lies in $N(\mathcal S)$. If
$E'E^\dagger\in\mathcal S$ the two act identically on the code space (a
harmless \emph{degeneracy}); if $E'E^\dagger\in N(\mathcal
S)\setminus\mathcal S$ they differ by a nontrivial logical operation and
are genuinely dangerous. Hence the dangerous errors are exactly the
low-weight logical operators, and a code of distance $d$ --- the minimum
such weight --- distinguishes (and so corrects) all errors of weight up
to $\lfloor(d-1)/2\rfloor$, since two correctable errors then differ by
weight $<d$ and cannot both be logical. With $n-k$ generators there are
$2^{n-k}$ syndromes available to label the correctable errors.
\end{proof}

\begin{statement}[\textbf{Knill--Laflamme conditions}]
\label{stmt:knill-laflamme}
Equivalently, a code with orthonormal logical basis $\{\ket{i_L}\}$
corrects an error set $\{E_a\}$ if and only if
$\bra{i_L}E_a^\dagger E_b\ket{j_L}=C_{ab}\,\delta_{ij}$ for a Hermitian
matrix $C_{ab}$ independent of the logical indices. The diagonal form in
$i,j$ says errors never leak information about the logical state; the
single matrix $C_{ab}$ says different errors map the code space to
distinguishable (or identical) subspaces.
\end{statement}

\begin{example}[The repetition code in the formalism]
\label{ex:rep-stabiliser}
The bit-flip code is the $[[3,1,1]]$ stabiliser code with
$\mathcal S=\langle\hat\sigma_z^{(1)}\hat\sigma_z^{(2)},
\hat\sigma_z^{(2)}\hat\sigma_z^{(3)}\rangle$, logical operators
$\op{Z}_L=\hat\sigma_z^{(1)}$ and
$\op{X}_L=\hat\sigma_x^{(1)}\hat\sigma_x^{(2)}\hat\sigma_x^{(3)}$. Its
distance is $d=1$ because the weight-one $\hat\sigma_z^{(1)}$ is a
logical operator --- the precise statement that phase errors go
uncorrected. As a code against bit flips \emph{only}, however, its
$\hat\sigma_x$-distance is $3$ and it corrects one flip, recovering the
$\lfloor(3-1)/2\rfloor=1$ count. Reaching a true distance-$3$ quantum
code, correcting \emph{any} single-qubit error, requires more qubits.
\end{example}

\paragraph{Transition.}
The formalism makes the goal concrete: find commuting Pauli generators
whose smallest nontrivial logical operator has weight $3$, so that every
single-qubit $\hat\sigma_x$, $\hat\sigma_y$, $\hat\sigma_z$ is detected
and located. Two classic solutions --- one by nesting repetition codes,
one from a classical Hamming code --- are the CSS codes.

% --------------------------------------------------------------
\section{CSS codes: the Shor and Steane codes}
\label{sec:css}

\begin{phenomenon}
A code that corrects both bit and phase flips can be assembled from
two classical codes acting in the two conjugate bases: one set of checks
($\hat\sigma_z$-type) catches bit flips, an independent set
($\hat\sigma_x$-type) catches phase flips. When the two classical codes
fit together compatibly --- so that the $\hat\sigma_x$- and
$\hat\sigma_z$-checks all commute --- the result is a
Calderbank--Shor--Steane (CSS) code. Shor's nine-qubit code builds it by
brute-force nesting; Steane's seven-qubit code builds it elegantly from
a single Hamming code.
\end{phenomenon}

\begin{statement}[\textbf{CSS codes, Shor, and Steane}~\cite{shor1995,steane1996,calderbank1996}]
\label{stmt:css}
The \emph{Shor code} $[[9,1,3]]$ encodes
\begin{equation}
\ket{0_L}=\tfrac{1}{2\sqrt2}(\ket{000}+\ket{111})^{\otimes3},
\qquad
\ket{1_L}=\tfrac{1}{2\sqrt2}(\ket{000}-\ket{111})^{\otimes3},
\end{equation}
with eight stabiliser generators: six $\hat\sigma_z$-type checks within
the three blocks and two $\hat\sigma_x$-type checks across them,
\begin{equation}
\begin{aligned}
&\hat\sigma_z^{(1)}\hat\sigma_z^{(2)},\;
\hat\sigma_z^{(2)}\hat\sigma_z^{(3)},\;
\hat\sigma_z^{(4)}\hat\sigma_z^{(5)},\;
\hat\sigma_z^{(5)}\hat\sigma_z^{(6)},\;
\hat\sigma_z^{(7)}\hat\sigma_z^{(8)},\;
\hat\sigma_z^{(8)}\hat\sigma_z^{(9)},\\
&\hat\sigma_x^{(1)}\hat\sigma_x^{(2)}\hat\sigma_x^{(3)}
\hat\sigma_x^{(4)}\hat\sigma_x^{(5)}\hat\sigma_x^{(6)},\quad
\hat\sigma_x^{(4)}\hat\sigma_x^{(5)}\hat\sigma_x^{(6)}
\hat\sigma_x^{(7)}\hat\sigma_x^{(8)}\hat\sigma_x^{(9)}.
\end{aligned}
\end{equation}
The \emph{Steane code} $[[7,1,3]]$ is the CSS code built from the
$[7,4,3]$ Hamming code with parity-check matrix
\begin{equation}
\label{eq:hamming}
H=\begin{pmatrix}
0&0&0&1&1&1&1\\
0&1&1&0&0&1&1\\
1&0&1&0&1&0&1
\end{pmatrix},
\end{equation}
whose three rows give three $\hat\sigma_z$-type and three
$\hat\sigma_x$-type generators (a $\hat\sigma_z$ or $\hat\sigma_x$ on
each qubit where the row is $1$). Both codes have distance $3$ and
correct any single-qubit error.
\end{statement}

\begin{proof}[Derivation]
\emph{Shor.} Read the code as a phase-flip code of three blocks, each
block itself a bit-flip code. The six intra-block
$\hat\sigma_z\hat\sigma_z$ checks locate a bit flip inside whichever
block it hits, exactly as in \cref{stmt:repetition}: an
$\hat\sigma_x^{(5)}$ error makes qubits $4,5$ and $5,6$ disagree, firing
$\hat\sigma_z^{(4)}\hat\sigma_z^{(5)}$ and
$\hat\sigma_z^{(5)}\hat\sigma_z^{(6)}$ and pinpointing qubit $5$. The two
inter-block $\hat\sigma_x^{\cdots}$ checks compare the \emph{signs} of
adjacent blocks: a phase flip $\hat\sigma_z^{(j)}$ anywhere in a block
flips that block's sign $(\ket{000}+\ket{111})\to(\ket{000}-\ket{111})$,
which the weight-six $\hat\sigma_x$ checks detect by comparing block~1
with block~2 and block~2 with block~3. A single $\hat\sigma_z$ in block
$2$ fires both $\hat\sigma_x$-checks; recovery applies $\hat\sigma_z$ to
any qubit of that block, the choice being a stabiliser degeneracy. A
$\hat\sigma_y$ error trips both a $\hat\sigma_z$- and an
$\hat\sigma_x$-check and is corrected as the product. \emph{Steane.}
Because the Hamming code contains its dual, the $\hat\sigma_x$- and
$\hat\sigma_z$-generators built from the same rows of $H$ overlap in an
even number of qubits and hence commute, so $\mathcal S$ is abelian. For
an $\hat\sigma_x^{(j)}$ error the $\hat\sigma_z$-syndrome is the $j$th
column of $H$ --- the binary address of $j$ --- so the syndrome
\emph{is} the error location. For instance an error on qubit $6$ gives
column $(1,1,0)^{\!\top}$, the binary $6$. Phase errors are located
identically by the $\hat\sigma_x$-checks. The minimum-weight logical
operators are the transversal
$\op{X}_L=\hat\sigma_x^{\otimes7}$ and
$\op{Z}_L=\hat\sigma_z^{\otimes7}$, giving distance $3$.
\end{proof}

\begin{remark}
\textbf{Syndrome extraction in practice.} Each generator is measured
with one ancilla. For a $\hat\sigma_z$-type check
$\hat\sigma_z^{(a)}\hat\sigma_z^{(b)}$, prepare an ancilla in $\ket{0}$,
apply $\mathrm{CNOT}$ from qubits $a$ and $b$ onto it, and measure it:
the ancilla reads the parity $\hat\sigma_z^{(a)}\hat\sigma_z^{(b)}$
without revealing either bit. An $\hat\sigma_x$-type check is the same
circuit conjugated by Hadamards on the ancilla (so it measures
$\hat\sigma_x$ parity). The full code runs all $n-k$ such gadgets each
cycle; fault-tolerant designs repeat and cross-check them so that an
error in the syndrome circuit itself does not masquerade as a data
error.
\end{remark}

\begin{application}[Transversal Cliffords and the return of magic states]
\label{app:transversal}
The Steane code's CSS structure makes the Clifford gates
$\{\mathrm{H},\mathrm{S},\mathrm{CNOT}\}$ \emph{transversal}: applying
the physical gate to all seven qubits of a block (or qubit-wise between
two blocks for CNOT) enacts the logical gate, so no single physical
fault spreads within a block. This is the cleanest possible fault
tolerance --- but by the Eastin--Knill theorem (\cref{app:magic-cost})
it cannot extend to a universal set. The non-Clifford $\mathrm{T}$ is not
transversal in the Steane code and must be supplied through
magic-state distillation (\cref{stmt:magic-distillation}), confirming
from the coding side the dichotomy isolated in the previous chapter:
Cliffords are cheap and protected, the $\mathrm{T}$ gate is expensive and
imported.
\end{application}

\paragraph{Transition.}
Shor and Steane prove fault tolerance is possible, but their stabilisers
are geometrically awkward --- weight-six operators spanning distant
qubits. Hardware wants checks that are local on a two-dimensional chip.
The surface code delivers exactly that, and with it the highest known
error threshold.

% --------------------------------------------------------------
\section{The surface code}
\label{sec:surface-code}

\begin{phenomenon}
A superconducting processor (\cref{ch:experimental-basics}) couples each
qubit only to its nearest neighbours on a planar chip. A practical code
must therefore use checks that are geometrically local and low weight.
The surface code does this with a single repeating pattern of weight-four
parity checks on a two-dimensional lattice, tolerates an error rate near
$1\%$, and is the architecture behind the recent below-threshold
demonstrations in cQED. Its price is overhead --- hundreds to thousands
of physical qubits per logical qubit --- but its locality is what makes
it buildable.
\end{phenomenon}

\begin{statement}[\textbf{The surface code}~\cite{kitaev2003,fowler2012}]
\label{stmt:surface-code}
Place data qubits on the edges of a two-dimensional square lattice.
Define two families of stabilisers: a \emph{star} (vertex) operator at
each vertex $v$,
\begin{equation}
\op A_v=\prod_{e\,\ni\,v}\hat\sigma_x^{(e)},
\qquad
\op B_p=\prod_{e\,\in\,\partial p}\hat\sigma_z^{(e)},
\end{equation}
the product of $\hat\sigma_x$ over the four edges meeting at $v$, and a
\emph{plaquette} operator $\op B_p$, the product of $\hat\sigma_z$ over
the four edges bounding each face $p$. All stars and plaquettes commute,
so the code space is their joint $+1$ eigenspace. Logical operators are
$\hat\sigma_x$- and $\hat\sigma_z$-strings running across the patch from
one boundary to the opposite one; on a distance-$d$ patch (linear size
$d$) they have weight $d$, encoding one logical qubit that corrects
$\lfloor(d-1)/2\rfloor$ errors. Below a threshold physical error rate
$p_{\rm th}\!\sim\!1\%$, the logical error rate falls as
\begin{equation}
\label{eq:threshold}
p_L\sim\Bigl(\frac{p}{p_{\rm th}}\Bigr)^{\!\lfloor d/2\rfloor},
\end{equation}
the threshold theorem of fault-tolerant computation.
\end{statement}

\begin{proof}[Derivation]
\emph{Commutation.} A star $\op A_v$ and a plaquette $\op B_p$ share
either no edges or exactly two (when $v$ sits on the boundary of $p$).
Where they share an edge, $\hat\sigma_x$ and $\hat\sigma_z$ anticommute;
with an even number of shared edges the anticommutations cancel and
$[\op A_v,\op B_p]=0$. Stars commute with stars and plaquettes with
plaquettes trivially (same Pauli type). The stabiliser group is
therefore abelian. \emph{Errors as strings.} A single $\hat\sigma_x$
error on an edge anticommutes with the two plaquettes that border that
edge, lighting up a pair of $\hat\sigma_z$-syndrome defects; a connected
chain of $\hat\sigma_x$ errors creates defects only at its two endpoints,
because interior plaquettes see two errors and stay $+1$. The decoder
sees the endpoint defects and pairs them with the shortest connecting
chain (minimum-weight perfect matching); it fails only if it closes the
chain the wrong way around the lattice. \emph{Distance.} A logical
$\op{X}_L$ is an unbroken $\hat\sigma_x$-string from boundary to
boundary: it commutes with every check (each interior vertex it crosses
shares two edges with it) yet is not a product of stars, so it is a
genuine logical operation. The shortest such string has weight $d$, the
patch size, giving the distance and hence \cref{eq:threshold}: a logical
fault needs an error chain spanning the lattice, of length $\sim d/2$
after matching, whose probability is exponentially small in $d$ once
$p<p_{\rm th}$.
\end{proof}

\begin{example}[A single error and its syndrome]
\label{ex:surface-syndrome}
Consider one $\hat\sigma_z$ error on an interior edge. It anticommutes
with the two \emph{star} operators at the edge's endpoints, so exactly
two adjacent $\hat\sigma_x$-syndrome bits flip while every plaquette
stays quiet --- an unmistakable signature pointing at the edge between
them. A neighbouring $\hat\sigma_z$ error sharing one of those vertices
would extend the chain, moving the defect outward rather than creating a
third one. This is why the decoder reasons about \emph{pairs} of defects
and the strings joining them, and why a distance-$3$ patch (correcting
any single error) is the smallest useful unit --- the building block of
the lattice-surgery and magic-state machinery that a full
fault-tolerant processor assembles.
\end{example}

\begin{application}[The surface code in cQED]
\label{app:surface-cqed}
The surface code maps directly onto a transmon lattice
(\cref{ch:experimental-basics}): data and ancilla qubits alternate on a
grid, each weight-four check is four CNOTs to a measurement ancilla read
out dispersively (\cref{stmt:dispersive-readout}), and the cycle repeats
every microsecond. Google's distance-$3$ and distance-$5$ logical
qubits~\cite{acharya2023} showed the logical error rate \emph{decreasing}
with code size --- the first experimental crossing of the threshold of
\cref{eq:threshold}, and the result the
sensing chapter (\cref{ch:coherent-multiqubit}) invokes when it argues
that error correction can in principle protect a Heisenberg-limited
probe.
\end{application}

\paragraph{Transition.}
The surface code spends many physical qubits to manufacture one good
logical qubit. A complementary philosophy asks whether the redundancy
can instead live in the infinite ladder of a single oscillator ---
trading qubit count for the rich state space of a superconducting cavity.

% --------------------------------------------------------------
\section{Bosonic codes in circuit QED}
\label{sec:bosonic-codes}

\begin{phenomenon}
A high-$Q$ superconducting cavity (\cref{ch:em-quantisation,ch:circuit-qed})
is a single mode with infinitely many levels and one dominant error:
single-photon loss at rate set by its $T_1$ (\cref{stmt:t1}). Rather than
stitch a logical qubit out of many two-level systems, a bosonic code
hides one logical qubit in the cavity's Fock ladder and corrects that
single dominant error using a transmon only to read a parity syndrome.
This \emph{hardware-efficient} approach has reached the milestone where a
logical qubit outlives its best physical component.
\end{phenomenon}

\begin{statement}[\textbf{Cat and GKP codes}~\cite{mirrahimi2014,gkp2001,michael2016}]
\label{stmt:bosonic-codes}
A \emph{cat code} encodes the logical qubit in superpositions of coherent
states (\cref{stmt:coherent}) of opposite phase,
\begin{equation}
\ket{\mathcal C^+_\alpha}\propto\ket{\alpha}+\ket{-\alpha}
\;(\text{even photon number}),
\qquad
\ket{\mathcal C^-_\alpha}\propto\ket{\alpha}-\ket{-\alpha}
\;(\text{odd}),
\end{equation}
so that photon loss $\op a$ toggles the photon-number parity
$\op P=e^{i\pi\op a^\dagger\op a}$: parity is the error syndrome, read
out without demolition through a dispersively coupled transmon
(\cref{stmt:dispersive-readout}). A \emph{GKP code} instead encodes the
qubit in grid states of a single mode, stabilised by the two commuting
displacements
\begin{equation}
\label{eq:gkp}
\op S_q=e^{\,i2\sqrt{\pi}\,\op q},
\qquad
\op S_p=e^{\,i2\sqrt{\pi}\,\op p},
\end{equation}
in dimensionless quadratures $[\op q,\op p]=i$ (\cref{stmt:quadratures}),
and corrects small displacements in both quadratures. Both codes have
demonstrated error-corrected logical qubits beyond the break-even point.
\end{statement}

\begin{proof}[Derivation]
\emph{Cat parity.} Expanding a coherent state in Fock states
(\cref{stmt:coherent}), $\ket{\pm\alpha}=e^{-|\alpha|^2/2}\sum_n
(\pm\alpha)^n/\sqrt{n!}\,\ket{n}$, the even combination
$\ket{\alpha}+\ket{-\alpha}$ keeps only even $n$ and the odd combination
only odd $n$; hence $\op P\ket{\mathcal C^\pm_\alpha}
=\pm\ket{\mathcal C^\pm_\alpha}$. A single photon loss acts as the
annihilation operator, $\op a\ket{\alpha}=\alpha\ket{\alpha}$, so
$\op a(\ket{\alpha}\pm\ket{-\alpha})=\alpha(\ket{\alpha}\mp\ket{-\alpha})$
--- it maps an even cat to an odd one and vice versa, flipping the
parity. Crucially it does \emph{not} mix the two logical states beyond
this known flip: monitoring $\op P$ between jumps tells us how many
photons have leaked, and the logical information (which superposition)
is preserved up to a tracked, correctable parity frame. The transmon
measures $\op P$ as a QND parity check because the dispersive coupling
$\propto\hat\sigma_z\,\op a^\dagger\op a$ rotates the transmon by an
angle proportional to photon number, distinguishing even from odd
without revealing $n$ itself. \emph{GKP stabilisers.} Using the Weyl
relation $e^{ia\op q}e^{ib\op p}=e^{-iab}e^{ib\op p}e^{ia\op q}$, the two
displacements in \cref{eq:gkp} satisfy
$\op S_q\op S_p=e^{-i(2\sqrt\pi)^2}\op S_p\op S_q=e^{-i4\pi}\op S_p\op S_q
=\op S_p\op S_q$ and so commute, defining a valid code space: the
grid of points spaced by $2\sqrt\pi$ in each quadrature. The logical
operators are the half-period translations
$\op{Z}_L=e^{i\sqrt\pi\,\op q}$ and $\op{X}_L=e^{i\sqrt\pi\,\op p}$, which
commute with the stabilisers (their phase is $e^{-i2\pi}=1$) but
anticommute with each other ($e^{-i\pi}=-1$), as a logical pair must. A
small displacement error shifts the grid by less than half a cell;
measuring the stabilisers modulo $2\sqrt\pi$ reveals the shift, and a
counter-displacement undoes it.
\end{proof}

\begin{application}[Beyond break-even]
\label{app:break-even}
Bosonic codes turn the cavity's single dominant loss channel into a
measurable, correctable event, and they have crossed the threshold that
defines useful QEC: an error-corrected logical qubit living longer than
the best uncorrected element of the same hardware. Photon-number-parity
tracking of a cat state in a 3D cavity reached the break-even point for
memory~\cite{ofek2016}, and real-time GKP correction has now surpassed
it, extending the logical lifetime beyond that of the bare
mode~\cite{sivak2023}. Because the encoding lives in one cavity read out
by one transmon, the hardware overhead is far smaller than a surface-code
patch --- at the cost of correcting a narrower class of errors. The two
philosophies, many-qubit surface codes and few-mode bosonic codes, are
the live alternatives for the logical qubits that a fault-tolerant
machine, and the long-lived sensing probes of \cref{ch:coherent-multiqubit},
will be built from.
\end{application}

% --------------------------------------------------------------
This chapter built quantum error correction from the ground up. The
obstacles that seemed to forbid it --- no-cloning, the continuum of
errors, measurement collapse --- were each turned aside: we encode by
entanglement rather than copying, digitise continuous errors by
measuring a syndrome, and design the parity checks to commute with the
logical operators so the encoded state is never read. The stabiliser
formalism organised every code as a set of commuting Pauli checks whose
low-weight logical operators set the distance, from the three-qubit
repetition code through the Shor and Steane CSS codes to the locally
checkable surface code with its $\sim\!1\%$ threshold, and finally the
cat and GKP codes that protect a logical qubit inside a single cavity.
Along the way the Clifford/magic dichotomy of the previous chapter
reappeared as the transversality of CSS Cliffords and the irreducible
cost of the $\mathrm{T}$ gate.

This completes Part~V. We now possess the full quantum-information
toolkit --- gates, protocols, and the means to protect logical states
against the decoherence of Part~II. Part~VI turns the same hardware and
the same coherence-engineering to a different end: instead of computing
with the qubit, it uses the qubit's exquisite sensitivity as a detector.
The dispersive readout, the long coherence times, and --- in the closing
chapter --- the entanglement and error-correction ideas of this part all
return, now in the service of sensing dark matter and single quanta.

\clearpage
\begin{chaptersummary}

\paragraph{Why QEC works.}
Encode without cloning (entangle one logical qubit across many physical
ones), measure without collapse (parity checks commuting with logical
operators), and digitise the continuum: correcting
$\{\hat\sigma_x,\hat\sigma_y,\hat\sigma_z\}$ on each qubit corrects any
error, because the syndrome measurement projects a continuous error onto
a definite Pauli branch.

\paragraph{Repetition and stabilisers.}
Bit-flip code $\ket{0_L}=\ket{000}$, $\ket{1_L}=\ket{111}$, stabilisers
$\hat\sigma_z^{(1)}\hat\sigma_z^{(2)},\hat\sigma_z^{(2)}\hat\sigma_z^{(3)}$;
the two-bit syndrome locates one flip. Generally an $[[n,k,d]]$ code is
the $+1$ eigenspace of $n-k$ commuting Paulis; an error is flagged when
it anticommutes with a generator, logical operators are the nontrivial
Paulis commuting with all of them, and distance $d$ (their minimum
weight) corrects $\lfloor(d-1)/2\rfloor$ errors.

\paragraph{CSS codes.}
Separate $\hat\sigma_x$- and $\hat\sigma_z$-checks correct phase and bit
flips. Shor $[[9,1,3]]$ nests a bit-flip code inside a phase-flip code;
Steane $[[7,1,3]]$ is built from the $[7,4,3]$ Hamming code, with the
$\hat\sigma_x$-syndrome equal to the binary address of the flipped
qubit and transversal Clifford gates --- but no transversal $\mathrm{T}$
(Eastin--Knill), so magic states return.

\paragraph{Surface code.}
Local weight-four checks $\op A_v=\prod\hat\sigma_x$ (stars),
$\op B_p=\prod\hat\sigma_z$ (plaquettes) on a 2D lattice; they commute
(share $0$ or $2$ edges). Errors are strings flagged at their endpoints;
logical operators span the patch with weight $d$. Below
$p_{\rm th}\!\sim\!1\%$, $p_L\sim(p/p_{\rm th})^{\lfloor d/2\rfloor}$.

\paragraph{Bosonic codes.}
Cat code: $\ket{\alpha}\pm\ket{-\alpha}$ with photon-number parity
$\op P=e^{i\pi\op a^\dagger\op a}$ as the loss syndrome, read QND by a
dispersively coupled transmon. GKP code: grid states stabilised by
displacements $e^{i2\sqrt\pi\op q},e^{i2\sqrt\pi\op p}$, correcting small
shifts. Both have surpassed break-even in cQED.

\end{chaptersummary}

% --------------------------------------------------------------
\section*{Exercises}
\addcontentsline{toc}{section}{Exercises}

\begin{exercise}[Reading the repetition-code syndrome]
The bit-flip code is measured and returns the syndrome
$(g_1,g_2)=(-1,-1)$. Which qubit was flipped, what recovery is applied,
and why does the measurement not reveal $\alpha,\beta$?
\begin{solution}
From the syndrome table in \cref{stmt:repetition}, $(-1,-1)$ means qubits
$1,2$ disagree and qubits $2,3$ disagree, so qubit $2$ is the odd one
out: a $\hat\sigma_x^{(2)}$ error. Recovery applies $\hat\sigma_x^{(2)}$.
The checks $g_1,g_2$ are products of $\hat\sigma_z$ and commute with both
logical operators, so they measure only \emph{relative} parities of the
three qubits, never the logical amplitudes --- the encoded state emerges
untouched.
\end{solution}
\end{exercise}

\begin{exercise}[Parity checks commute with the logicals]
Verify explicitly that $g_1=\hat\sigma_z^{(1)}\hat\sigma_z^{(2)}$
commutes with the logical operator
$\op{X}_L=\hat\sigma_x^{(1)}\hat\sigma_x^{(2)}\hat\sigma_x^{(3)}$.
\begin{solution}
$g_1$ and $\op{X}_L$ overlap on qubits $1$ and $2$. On qubit $1$,
$\hat\sigma_z^{(1)}$ and $\hat\sigma_x^{(1)}$ anticommute; on qubit $2$,
$\hat\sigma_z^{(2)}$ and $\hat\sigma_x^{(2)}$ anticommute; on qubit $3$,
$g_1$ acts as identity and there is no anticommutation. Two
anticommuting pairs give an overall factor $(-1)^2=+1$, so
$g_1\op{X}_L=\op{X}_L g_1$. The check therefore cannot disturb the
logical bit.
\end{solution}
\end{exercise}

\begin{exercise}[Digitising a coherent error]
A qubit of the bit-flip code suffers the coherent rotation
$\op R=\cos\theta\,\id-i\sin\theta\,\hat\sigma_x^{(1)}$. After syndrome
measurement, what are the possible outcomes and their probabilities, and
is the code word recovered in each case?
\begin{solution}
$\op R\ket{\psi_L}=\cos\theta\,\ket{\psi_L}
-i\sin\theta\,\hat\sigma_x^{(1)}\ket{\psi_L}$. Measuring $(g_1,g_2)$
projects onto the trivial syndrome with probability $\cos^2\theta$,
leaving $\ket{\psi_L}$ (nothing to do), or onto the
$\hat\sigma_x^{(1)}$ syndrome $(-1,+1)$ with probability $\sin^2\theta$,
leaving $\hat\sigma_x^{(1)}\ket{\psi_L}$, which is corrected by applying
$\hat\sigma_x^{(1)}$. In both branches the exact code word is restored:
the continuous parameter $\theta$ has been digitised
(\cref{stmt:digitisation}).
\end{solution}
\end{exercise}

\begin{exercise}[The phase-flip code]
Write the code words, stabilisers, and logical operators of the
three-qubit phase-flip code, and explain its relation to the bit-flip
code in one sentence.
\begin{solution}
Code words $\ket{0_L}=\ket{+\!+\!+}$, $\ket{1_L}=\ket{-\!-\!-}$;
stabilisers $\hat\sigma_x^{(1)}\hat\sigma_x^{(2)}$ and
$\hat\sigma_x^{(2)}\hat\sigma_x^{(3)}$; logical operators
$\op{X}_L=\hat\sigma_x^{(1)}$ and
$\op{Z}_L=\hat\sigma_z^{(1)}\hat\sigma_z^{(2)}\hat\sigma_z^{(3)}$. It is
the bit-flip code conjugated by a Hadamard on every qubit, which swaps
$\hat\sigma_x\leftrightarrow\hat\sigma_z$ (\cref{eq:clifford-conjugation})
and hence swaps the roles of bit and phase flips.
\end{solution}
\end{exercise}

\begin{exercise}[Shor-code syndromes]
In the Shor code, which stabilisers fire for (a) a bit flip
$\hat\sigma_x^{(5)}$ and (b) a phase flip $\hat\sigma_z^{(5)}$?
\begin{solution}
(a) $\hat\sigma_x^{(5)}$ anticommutes with the intra-block
$\hat\sigma_z$-checks that contain qubit $5$, namely
$\hat\sigma_z^{(4)}\hat\sigma_z^{(5)}$ and
$\hat\sigma_z^{(5)}\hat\sigma_z^{(6)}$; both fire and locate qubit $5$.
The $\hat\sigma_x$-checks commute with it and stay quiet. (b)
$\hat\sigma_z^{(5)}$ anticommutes with the two weight-six
$\hat\sigma_x$-checks, since qubit $5$ lies in both
$\hat\sigma_x^{(1)}\!\cdots\hat\sigma_x^{(6)}$ and
$\hat\sigma_x^{(4)}\!\cdots\hat\sigma_x^{(9)}$; both fire, flagging a
phase flip in block~2. The recovery $\hat\sigma_z$ on any block-2 qubit
works, as the choice differs only by a stabiliser.
\end{solution}
\end{exercise}

\begin{exercise}[The Steane syndrome is an address]
Using the Hamming matrix \cref{eq:hamming}, find the
$\hat\sigma_z$-syndrome of a bit-flip error on qubit $3$, and confirm it
equals the binary representation of $3$.
\begin{solution}
The syndrome of an $\hat\sigma_x^{(j)}$ error is column $j$ of $H$. Column
$3$ is $(0,1,1)^{\!\top}$: the top row has $0$ in position $3$, the
middle and bottom rows have $1$. Reading top-to-bottom as bits of value
$(4,2,1)$ gives $0\cdot4+1\cdot2+1\cdot1=3$, the binary address of the
flipped qubit. The three $\hat\sigma_z$-checks
$\hat\sigma_z^{(2)}\hat\sigma_z^{(3)}\hat\sigma_z^{(6)}\hat\sigma_z^{(7)}$
and $\hat\sigma_z^{(1)}\hat\sigma_z^{(3)}\hat\sigma_z^{(5)}\hat\sigma_z^{(7)}$
fire while the first stays quiet, directly pinpointing qubit $3$.
\end{solution}
\end{exercise}

\begin{exercise}[Counting syndromes]
The Steane code is $[[7,1,3]]$. How many stabiliser generators does it
have, how many distinct syndromes, and is that enough to label every
correctable single-qubit Pauli error?
\begin{solution}
There are $n-k=7-1=6$ generators, giving $2^6=64$ syndromes. The
correctable weight-one errors are $\hat\sigma_x,\hat\sigma_y,\hat\sigma_z$
on each of $7$ qubits, $3\times7=21$ nontrivial errors, plus the trivial
(no-error) syndrome. Since $21+1=22\le64$ and the code's distance $3$
guarantees these have distinct syndromes, the syndrome space is more than
sufficient --- a single-qubit error is always uniquely diagnosed.
\end{solution}
\end{exercise}

\begin{exercise}[Surface-code commutation]
Show that a star $\op A_v$ and a plaquette $\op B_p$ that share exactly
two edges commute.
\begin{solution}
On each shared edge, $\op A_v$ contributes $\hat\sigma_x$ and $\op B_p$
contributes $\hat\sigma_z$; these anticommute. On all other edges at
most one of the operators acts, contributing commuting factors. With two
shared edges there are two anticommuting pairs, so swapping the operators
incurs $(-1)^2=+1$ and $[\op A_v,\op B_p]=0$. (An odd overlap is
forbidden by the lattice geometry --- a vertex always meets a plaquette
boundary in $0$ or $2$ edges --- which is exactly why the surface-code
stabilisers form an abelian group.)
\end{solution}
\end{exercise}

\begin{exercise}[Distance and correctable errors on a patch]
A surface-code patch has linear size $d=5$. What is the minimum weight of
a logical operator, and how many arbitrary errors is the patch guaranteed
to correct?
\begin{solution}
A logical operator is a Pauli string crossing the patch from boundary to
boundary; the shortest such string has weight equal to the patch size,
$d=5$. The code therefore has distance $5$ and corrects any error of
weight up to $\lfloor(d-1)/2\rfloor=\lfloor4/2\rfloor=2$. Equivalently,
three well-placed errors can in principle form half of a logical string
and cause a logical fault, but any one or two are always corrected.
\end{solution}
\end{exercise}

\begin{exercise}[Loss flips a cat's parity]
Show that single-photon loss maps an even cat state to an odd one, and
explain why a dispersively coupled transmon can read this parity without
collapsing the logical superposition.
\begin{solution}
With $\ket{\mathcal C^\pm_\alpha}\propto\ket{\alpha}\pm\ket{-\alpha}$ and
$\op a\ket{\pm\alpha}=\pm\alpha\ket{\pm\alpha}$,
$\op a(\ket{\alpha}+\ket{-\alpha})=\alpha(\ket{\alpha}-\ket{-\alpha})$ ---
an even cat becomes an odd cat, i.e.\ the photon-number parity
$\op P=e^{i\pi\op a^\dagger\op a}$ flips from $+1$ to $-1$. The dispersive
interaction $\propto\hat\sigma_z\,\op a^\dagger\op a$
(\cref{stmt:dispersive-readout}) rotates the transmon by an angle
proportional to the photon number, so measuring the transmon distinguishes
even from odd $n$ \emph{without} resolving $n$ itself. Because the
measured observable is the parity, which commutes with the logical
operators of the cat code, the logical superposition (which cat) is
preserved while the loss event is detected.
\end{solution}
\end{exercise}
